%mac
%\begin{vidu}%Câu 183
%	Ta cần tác dụng một moment ngẫu lực $12\ N.m$ để làm quay bánh xe như hình. Xác định độ lớn lực tác dụng vào bánh xe ở hình a và hình b. Từ đó, hãy cho biết trường hợp nào sẽ có lợi hơn về lực.
%	\begin{center}
%		\includegraphics[width=8cm]{"IMG_01.-SBT - Ki 2 - Chuyen Latex/img-79"}
%	\end{center}
%\haidong{8}\loigiai{
%\begin{itemize}
%	\item 
%	\item 
%	\item 
%	\item 
%	\item 
%\end{itemize}}
%\end{vidu}
%\loai{Tính moment ngẫu lực}
\begin{vidu}%[1P3-8.2B][Loai1]
Hai lực song song, ngược chiều và có độ lớn bằng $5\ N$ đặt tại hai điểm $A$ và $B$, cánh tay đòn của ngẫu lực là $d = 10\ cm$ thì moment ngẫu lực bằng
\choice
{$50\ N$}
{\True $0,5\ N.m$}
{$0,5\ N$}
{$0,5\ J$}
\haidong{3}	\loigiai{
Ta có moment ngẫu lực: $M = F.d = 5.0,1 = 0.5$ N.m.
}
\end{vidu}
%\loai{Tính moment - hình tam giác}
\begin{vidu}%[1P3-8.2B][Loai1]
Một vật rắn phẳng, mỏng có dạng là một tam giác đều $ABC$, mỗi cạnh là $a = 20\ cm$. Người ta tác dụng vào vật một ngẫu lực nằm trong mặt phẳng của tam giác. Các lực có độ lớn là $8,0\ N$ đặt vào hai đỉnh $A$ và $B$. Tính moment của ngẫu lực khi các lực vuông góc với cạnh $AC$.
\choice
{$1,6\ N.m$}
{\True $0,8\ N.m$}
{$3,2\ N.m$}
{$1,38\ N.m$}
\haidong{5}\loigiai{
Ta có:	$M = F.d = F.\dfrac{a}{2} = 0,8$ N.m.
}
\end{vidu}
%\loai{Tính moment - hình vuông}
\begin{vidu}%[1P3-8.2B][Loai1]
Một vật rắn phẳng, mỏng có dạng hình vuông $ABCD$ cạnh là $a = 50\ cm$. Người ta tác dụng vào vật một ngẫu lực nằm trong mặt phẳng của hình vuông. Các lực có độ lớn $10\ N$ và đặt vào hai đỉnh $A$ và $C$. Moment của ngẫu lực trong trường hợp các lực vuông góc với $AC$ là
\choice
{$5\ N.m$}
{\True $5\sqrt 2\ N.m$}
{$500\ N.m$}
{$500\sqrt 2\ N.m$}
\haidong{6}	\loigiai{
Moment lực khi đó: $M = F.d = F.AC = F.a\sqrt 2 = 10.0,5\sqrt 2 = 5\sqrt 2\ N.m.$
}
\end{vidu}